\begin{abstract}
\thispagestyle{empty}
Este trabalho investiga a viabilidade do uso de reconhecimento facial baseado em \textit{embeddings}, representações vetoriais das características faciais, como alternativa aos cartões \textit{Near Field Communication} (\acs{NFC}) utilizados para identificação de usuários em ambientes institucionais.
Inicialmente, foi realizada uma validação conceitual e experimental da abordagem, utilizando modelos pré-treinados para extração de características faciais e métricas de similaridade no espaço vetorial. Essa etapa empregou uma base de imagens faciais sintéticas, adotada em conformidade com aspectos éticos e com a \acs{LGPD}, permitindo avaliar o desempenho por meio das métricas biométricas \acs{FAR} (\textit{False Acceptance Rate}) e \acs{FRR} (\textit{False Rejection Rate}).

Os resultados obtidos demonstraram boa separação entre identidades distintas e consistência entre amostras da mesma identidade, indicando a viabilidade técnica da solução. Com base nesses resultados, foi desenvolvido um protótipo funcional de reconhecimento facial em tempo real utilizando \textit{webcam}, armazenamento local de dados biométricos e interface gráfica para cadastro, gerenciamento e identificação de usuários. O sistema emprega os modelos \acs{MTCNN} para detecção facial e InceptionResnetV1 pré-treinado na base VGGFace2 para geração de \textit{embeddings}. Além disso, foram implementados mecanismos de cadastro multi-pose, avaliação da qualidade das capturas e votação temporal para aumentar a robustez do reconhecimento.

Os resultados obtidos mostram que a abordagem proposta é tecnicamente viável como alternativa aos cartões físicos, combinando evidências experimentais e validação prática por meio de um protótipo funcional capaz de realizar identificação biométrica de forma automatizada e sem contato físico.
\\[0.5em]
\textbf{Palavras-chave:} reconhecimento facial; biometria; \textit{embeddings}; visão computacional; controle de acesso; \acs{NFC}.
\end{abstract}
